\begin{textsample}{POS Dim 1 – vem\_ed\_08 – Score 29.00 – vem\_ed\_08/t033.txt}  \label{ex:f1_pos_064}
Olhar atento ao tempo Maira de Lourdes Rezende, doutora em Nanociências pela UFABC (2017), \textbf{é} pesquisadora e professora no Curso Superior de Tecnologia em Polímeros da Fatec Sorocaba.
No segundo semestre de 2020, participou da oficina Taller COIL Diseño, ofertada por um consórcio de universidades chilenas do qual a Universidad Tecnológica de Chile (INACAP) faz \textbf{parte}.
A \textbf{partir} de uma aproximação \textbf{proposta} pelas \textbf{coordenações} das duas instituições, Maira e Patricio Hernán Marabolí Albornoz, professor do curso de Engenharia de Minas da INACAP, desenvolveram na oficina o PCI"Soluções alternativas para a reciclagem dos resíduos \textbf{gerados} na indústria extrativa mineira a \textbf{partir} de um trabalho \textbf{colaborativo} e intercultural".
Desse PCI, participaram 12 estudantes brasileiros e 12 chilenos.
A \textbf{seguir}, o depoimento da professora sobre essa experiência: Elaboramos um projeto \textbf{interdisciplinar} (Engenharia de Minas e Tecnologia em Polímeros) em consonância com os Objetivos de Desenvolvimento \textbf{Sustentável} (ODS) da Organização das Nações Unidas (ONU).
Dessa \textbf{forma}, definimos a reciclagem de resíduos da indústria extrativa mineira como \textbf{foco} do projeto.
O planejamento se \textbf{baseou} no \textbf{cronograma} de aulas das duas IES, para desenvolver semanalmente as atividades (\textbf{síncronas} e assíncronas).
Vídeos, artigos científicos e distintas metodologias de ensino contribuíram para atingir o principal objetivo do PCI: proporcionar \textbf{vivência} intercultural e \textbf{promover} \textbf{desenvolvimento} \textbf{profissional} e pessoal dos estudantes.
O"quebra-gelo"(rompehielo, em espanhol) ocorreu nas duas primeiras semanas: na primeira, os alunos deveriam \textbf{escolher} 5 fotografias (suas ou não) \textbf{que} os \textbf{representassem}.
Nessa ocasião, os grupos se reuniram para conhecer os colegas e \textbf{criar} vínculos.
Na segunda atividade, aplicamos um questionário para \textbf{identificar} a"orientação cultural"dos estudantes, \textbf{que} \textbf{compartilharam} suas respostas entre os membros do grupo e elaboraram um vídeo em \textbf{que} \textbf{cada} integrante abordou um tópico do questionário.
A participação \textbf{foi} \textbf{ativa}, excelente.
O principal desafio, logo no \textbf{início} do PCI, \textbf{foi} a comunicação entre os alunos e a organização do tempo para entregar as \textbf{tarefas}.
Ferramentas como o WhatsApp facilitam a \textbf{interação}.
Além disso, realizar reuniões periódicas para esclarecer dúvidas, bem como o alinhamento e o engajamento dos docentes, \textbf{são} extremamente eficientes para incentivar o cumprimento das \textbf{tarefas}.
A avaliação ocorreu em \textbf{todas} as etapas, com elaboração de textos, apresentações ou vídeos pelas equipes.
Ao fim do PCI, os grupos relataram em vídeo a experiência, suas dificuldades e êxitos.
Para realizar um PCI deve-se considerar \textbf{idioma} da instituição parceira, objetivos de aprendizagem e detalhar \textbf{cronograma} e metodologias.
Aconselho àqueles professores \textbf{que} se interessem em realizar um PCI, \textbf{que} o façam.
A experiência adquirida, tanto no âmbito pessoal \textbf{quanto} no \textbf{profissional}, \textbf{é} incrível.

% matched lemmas: ativo, basear, cada, colaborativo, compartilhar, coordenação, criar, cronograma, desenvolvimento, escolher, foco, forma, gerar, identificar, idioma, interação, interdisciplinar, início, parte, partir, profissional, promover, propor, quanto, que, representar, seguir, ser, sustentável, síncrono, tarefa, todo, vivência
\end{textsample}
